\documentclass[11pt]{article}
\usepackage[margin=1in]{geometry}
\usepackage{amsmath,amssymb,amsthm,hyperref}
\newtheorem{theorem}{Theorem}
\title{Conjecture 00000001715: the displayed rotation multipliers do not decay}
\author{GitHub: gaochengzhecpu}
\date{}
\begin{document}
\maketitle
\begin{abstract}
The conjecture explicitly describes the coefficients
$c_k=(1-e(k\alpha))^{-1}$ as a decaying sequence. For every irrational
rotation parameter $\alpha$ and every positive integer $k$, these coefficients
are well defined and satisfy $|c_k|\geq 1/2$. Thus the stated decay assertion
is false. This disproves that necessary component of the conjecture; it does
not assert that every possible local limit theorem for rotations is false.
\end{abstract}

\section*{Statement and interpretation}
Both language versions of the source explicitly call
$c_k=(1-e(k\alpha))^{-1}$ a decaying sequence: the English says
``through the decaying sequence'', and the Chinese uses the corresponding
phrase for a decaying sequence. We use the standard notation
\[
 e(t)=\exp(2\pi i t),\qquad
 c_k(\alpha)=\frac{1}{1-\exp(2\pi i k\alpha)},\quad k\geq1.
\]
We restrict to irrational $\alpha$, the usual setting for irrational
rotations. Here decay has its ordinary meaning $c_k\to0$ as $k\to\infty$;
the argument also rules out convergence of the magnitudes to zero and any
positive power decay bound.

\begin{theorem}
For every irrational real number $\alpha$ and every integer $k\geq1$,
$c_k(\alpha)$ is well defined and $|c_k(\alpha)|\geq 1/2$.
Consequently $c_k(\alpha)$ does not tend to zero. In particular,
$\alpha=\sqrt2$ is an explicit counterexample to the asserted decay.
\end{theorem}
\begin{proof}
If $\exp(2\pi i k\alpha)=1$, the usual characterization of the kernel
of the complex exponential gives $k\alpha\in\mathbb Z$. Since $k$ is
a nonzero integer, this would make $\alpha$ rational. Therefore the
denominator is nonzero. The complex exponential on the imaginary axis
has modulus one, so the triangle inequality gives
\[
 0<|1-\exp(2\pi i k\alpha)|
 \leq |1|+|\exp(2\pi i k\alpha)|=2.
\]
Taking reciprocals yields $|c_k(\alpha)|\geq1/2$. If $c_k(\alpha)$
converged to zero, continuity of the norm would imply
$|c_k(\alpha)|\to0$, contradicting this uniform lower bound.
The standard irrationality of $\sqrt2$ supplies the stated explicit
parameter.
\end{proof}

\section*{Scope of the disproof}
The source combines an LLT correction claim with this explicit decay
description. A conjunction with the false decay assertion cannot hold
as written. We do not claim to settle a reformulation that omits decay
of $c_k$. In particular, multiplying these nondecaying coefficients by
Fourier coefficients of a suitably smooth function can produce decaying
products. Our result does not rule that out and does not analyze such
products, the LLT itself, or the source's additional condition about
low frequencies. It concerns the displayed sequence $c_k$ itself, exactly
as identified in both language versions. No assertion of unboundedness
is needed or made.

\section*{Formal verification and scope}
The accompanying Lean 4 file defines the actual complex exponential
$e(t)$ and the actual reciprocal coefficients. Its index $n\in\mathbb N$
corresponds to $k=n+1$, so it includes every positive frequency and
excludes the undefined zero frequency. It proves nonvanishing of the
denominators from Mathlib's \texttt{Irrational} predicate and the standard
complex exponential kernel theorem, proves the uniform norm bound, and
proves failure of convergence in the standard complex topology using
\texttt{Filter.Tendsto}. It also proves the counterexample at $\sqrt2$.

The universally quantified proposition named \texttt{DecayAssertion}
states decay for every irrational real rotation parameter. The final
theorem is its negation:
\[
\texttt{Conjecture1715.disprove\_decay\_assertion}.
\]
This proposition formalizes the necessary decay component, not an
invented definition of the entire, otherwise underspecified LLT claim.
The other principal theorems give the stronger result that decay fails
for every irrational parameter. All mathematics in the displayed theorem
is formalized; the source interpretation and logical relevance to its
conjunction are explained here. The axiom audit lists only
\texttt{propext}, \texttt{Classical.choice}, and \texttt{Quot.sound}.

\section*{Source and reproduction}
The original bilingual statement is included unchanged as
\texttt{SOURCE.md}, fetched from upstream commit
\begin{center}\small
\texttt{b355e42906e6eaf52bc7e9d005e0ef02d8aa7400}.
\end{center}
Its public source is
\href{https://github.com/The-Last-Math-Competition/The-Last-Math-Competition/blob/b355e42906e6eaf52bc7e9d005e0ef02d8aa7400/conjectures/00000001715.md}{conjecture 00000001715}.
The project uses Lean 4.19.0 and Mathlib commit
\begin{center}\small
\texttt{c44e0c8ee63ca166450922a373c7409c5d26b00b}.
\end{center}
From the \texttt{lean} directory, run \texttt{lake build}, then
\begin{center}
\texttt{lake env lean -DwarningAsError=true Main.lean}.
\end{center}
The main file ends with axiom-printing commands for the four principal
results. This draft was developed by a delegated AI agent with a
self-review; the coordinating agent reviews it separately.
\end{document}
